\section{Results}
\label{sec:results}

\subsection{The evaluator path moves the headline number more than any fault}
\label{sec:res-fallback}

Table~\ref{tab:baseline} stratifies the frozen control baseline by evaluator
path. On the native path the agent succeeds in \BaseNativeRate{} of
\BaseNativeN{} trials; on the compatibility fallback it succeeds in
\BaseFallbackRate{} of \BaseFallbackN{}. Every one of the fallback failures is
structural --- the fallback's null-retrieval branch is unreachable for these
tasks --- not a property of the agent. Pooling the two paths yields
\BasePooledRate{}, and that pooled number is what a pipeline without
stratification would report as ``the baseline.''

This is a \BaselineSwing{}~pp artefact of instrumentation, and it is larger than
any fault effect measured anywhere in this paper. Two rules follow, and we apply
them throughout: never pool the paths, and never compare an arm scored on one
path against an arm scored on the other.

% GENERATED by scripts/paper/make_tables.py -- do not edit by hand.
\begin{table}[t]
\centering
\small
\caption{Frozen 48-trial ReAct baseline, stratified by evaluator path. Native and fallback rows are \emph{not} comparable and must never be pooled; the pooled row is reported for completeness only.}
\label{tab:baseline}
\begin{adjustbox}{max width=\linewidth}
\begin{tabular}{lrrrrr}
\toprule
Stratum & Trials & Success & Success rate & Mean steps & Mean time (s) \\
\midrule
All trials & 48 & 25 & 52.1\% & 6.40 & 122.9 \\
Native evaluator & 20 & 17 & 85.0\% & 6.30 & 91.0 \\
Fallback evaluator & 28 & 8 & 28.6\% & 6.46 & 145.7 \\
\bottomrule
\end{tabular}
\end{adjustbox}
\end{table}


\subsection{No fault has a detectable effect on official success}
\label{sec:res-main}

Table~\ref{tab:official_native} gives the primary result on native-evaluator
pairs, and Figure~\ref{fig:forest} in \appref{app:figs} plots it. Reading the
point estimates alone, the picture is mixed and slightly
absurd: DOM omission is \ResWebDomMissingDeltaNat{}~pp, parameter errors
\ResAgentParamErrorDeltaNat{}~pp, but popups are
\emph{\ResWebPopupBlockDeltaNat{}~pp} and timeouts
\ResWebTimeoutDeltaNat{}~pp --- the fault arms doing \emph{better} than control.
Those signs are not evidence that popups help. They are what sampling noise looks
like at this sample size, and the exact McNemar tests say so: no fault reaches
significance (\ResAgentParamErrorPNat{} at the smallest, for parameter errors),
and after Holm correction every adjusted $p$ is exactly $1.000$.

The bootstrap intervals tell the same story. The interval for parameter errors,
the largest point estimate, spans
\ResAgentParamErrorDeltaNat{}~pp $\in$ $[-29.0, +0.0]$~pp --- a range that
includes both a substantial degradation and no effect at all. Reporting
``parameter errors do not significantly hurt the agent'' would be reporting the
width of that interval, not a property of the agent.

% GENERATED by scripts/paper/make_tables.py -- do not edit by hand.
\begin{table}[t]
\centering
\small
\caption{Primary fault result on \textbf{native-evaluator pairs only}. This is the table the paper relies on; the pooled version is contaminated by a fallback adapter that fails by construction. $n$ is below 40 because a pair is dropped when either arm was scored by the fallback adapter.}
\label{tab:official_native}
\begin{adjustbox}{max width=\linewidth}
\begin{tabular}{lrrrrrrrrrrrr}
\toprule
Fault & $n$ & Disc. & Control & Fault & Control\textsuperscript{\S} & Fault\textsuperscript{\S} & $\Delta$ (pp) & 95\% CI (pp) & $p_{\text{exact}}$ & $p_{\text{mid}}$ & $p_{\text{Holm}}$ & OR \\
\midrule
DOM omission & 30 & 10 & 83.3\% & 76.7\% & 84.4\% & 78.8\% & $-$6.7 & [$-$26.7, $+$13.3] & 0.754 & 0.549 & 1.000 & 0.69 \\
Popup block & 31 & 9 & 67.7\% & 77.4\% & 69.7\% & 77.4\% & $+$9.7 & [$-$9.7, $+$29.0] & 0.508 & 0.344 & 1.000 & 1.86 \\
HTTP error & 31 & 8 & 71.0\% & 71.0\% & 71.9\% & 73.5\% & $+$0.0 & [$-$19.4, $+$16.1] & 1.000 & 1.000 & 1.000 & 1.00 \\
Timeout & 31 & 3 & 77.4\% & 80.6\% & 78.8\% & 81.2\% & $+$3.2 & [$-$6.5, $+$12.9] & 1.000 & 0.625 & 1.000 & 1.67 \\
Parameter error & 31 & 6 & 93.5\% & 80.6\% & 93.9\% & 80.6\% & $-$12.9 & [$-$29.0, $+$0.0] & 0.219 & 0.125 & 1.000 & 0.27 \\
\textbf{All faults} & 154 & 36 & 78.6\% & 77.3\% & 79.8\% & 78.3\% & $-$1.3 & [$-$9.1, $+$6.5] & 0.868 & 0.743 & 1.000 & 0.90 \\
\bottomrule
\end{tabular}
\end{adjustbox}
\vspace{2pt}
\parbox{\linewidth}{\footnotesize $\Delta$ is the paired risk difference (fault $-$ control); the interval is a percentile bootstrap over pair units and is descriptive --- the exact McNemar $p$ is the primary inferential quantity. OR is the Haldane-corrected matched odds ratio. \textsuperscript{\S}Per-arm \emph{marginal} native rate, whose denominator is every native trial in that arm and therefore a \emph{different subset} from the pair-restricted rates in the Control/Fault columns. The two are not interchangeable: the paired rates are the ones $\Delta$ and the McNemar test refer to.}
\end{table}


For completeness, Table~\ref{tab:official_all} in \appref{app:pooled} reports the
same analysis with the evaluator paths pooled. It is not a robustness result and
we do not use it as one; we include it to show how much the conclusion moves when
the fallback contamination is left in.

\subsection{The design could not have detected a moderate effect}
\label{sec:res-power}

A null on its own is uninformative; what makes these nulls interpretable is the
power analysis. With \ResWebDomMissingNNat--\ResWebTimeoutNNat{} pairs per fault
and an observed discordance of \PiDNative{}, the paired design cannot reach
\Power{} power at \emph{any} effect size it can express --- the minimum detectable
effect is undefined, not merely large. At the observed discordance, detecting a
$10$~pp difference would require on the order of \PairsForTenPPObserved{} pairs
per fault, roughly \emph{six times} the \ResWebTimeoutNNat{} we have. Even a
counterfactual $n=40$ design reaches only \MdePairedNForty{}~pp. The full
power table, including the sample sizes implied by other discordances, is
Table~\ref{tab:power} in \appref{app:stats}.

The methodological consequence is that this family of null results is the default
outcome at realistic WebArena sample sizes, and should be reported as an
underpowered design rather than as evidence of robustness.

\subsection{The same faults are sharply visible at the process level}
\label{sec:res-behavior}

Table~\ref{tab:behavior} turns to the behavioural matrix, where the same five
faults are applied to \MatrixTasks{} tasks over \TrialsPerCell{} seeds, giving
\MatrixPairs{} pairs with step and timing instrumentation. Here the effects are
unambiguous, and they line up with mechanism rather than with layer:

\begin{itemize}
  \item \textbf{HTTP errors} add \BehWebHttpErrorStepDelta{} ReAct steps and
  \BehWebHttpErrorTimeDelta{}~s per pair (both $p<0.001$). An agent that
  encounters a server error re-navigates and retries; the cost is real and
  consistent.
  \item \textbf{Parameter errors} add \BehAgentParamErrorStepDelta{} steps
  ($p<0.001$): a corrupted argument costs a wasted action, the agent re-issues it,
  and the task usually still completes --- which is exactly why the success-rate
  effect is hard to see.
  \item \textbf{DOM omission} and \textbf{popup blocking} make the agent
  \emph{faster}, not slower (\BehWebDomMissingTimeDelta{}~s and
  \BehWebPopupBlockTimeDelta{}~s; the latter at $p=\BehWebPopupBlockTimeP{}$).
  We read this as fault-induced early termination rather than as a benefit: an
  agent that cannot see the page it needs gives up sooner.
\end{itemize}

Two cautions are built into how we report this. First, these rows carry no
official verdict --- no HAR was retained for this batch --- so ``cost'' and
``success'' in this paper come from different experiments over different task
sets, and we never combine them into a statement about cost per successful task.
Second, absolute cross-arm timings are not interpretable: the five per-fault
control arms are the same configuration, yet their mean wall-clock spans
\NoiseFloorLo--\NoiseFloorHi{}~s, a factor of \NoiseFloorRatio{}. Only
\emph{within-pair} deltas, which share the machine state, are reported.

% GENERATED by scripts/paper/make_tables.py -- do not edit by hand.
\begin{table}[t]
\centering
\small
\caption{Paired behavioral cost from the behavioural matrix (\MatrixPairs{} pairs, each logged twice). \textbf{Behavior-only}: no HAR was retained, so no row carries an official success verdict. Deltas are within-pair (fault $-$ control); $p$ is the two-sided Wilcoxon signed-rank test.}
\label{tab:behavior}
\begin{adjustbox}{max width=\linewidth}
\begin{tabular}{lrrrrrrr}
\toprule
Fault & $n$ pairs & $\Delta$ steps (mean) & $\Delta$ steps (HL) & $p$ & $\Delta$ time (s, mean) & $\Delta$ time (s, HL) & $p$ \\
\midrule
DOM omission & 80 & $-$0.10 & $+$0.0 & 0.956 & $-$66.1 & $-$46.0 & 0.068 \\
Popup block & 80 & $-$0.44 & $+$0.0 & 0.509 & $-$27.5 & $-$21.8 & 0.004 \\
HTTP error & 80 & $+$4.22 & $+$4.0 & $<$0.001 & $+$149.9 & $+$132.3 & $<$0.001 \\
Timeout & 80 & $-$0.72 & $-$0.5 & 0.021 & $-$13.2 & $-$8.9 & 0.316 \\
Parameter error & 80 & $+$2.30 & $+$1.5 & $<$0.001 & $+$20.6 & $+$7.1 & 0.460 \\
\bottomrule
\end{tabular}
\end{adjustbox}
\end{table}


\subsection{The architecture pilot cannot rank architectures}
\label{sec:res-arch}

Table~\ref{tab:architecture} in \appref{app:arch} gives the only architecture
comparison available. It is a pilot: $n=\ArchCellN$ per cell, three tasks, and ---
we established by inspection --- the committed runner does not reproduce these
rows, so the artifact and the script that supposedly produced it disagree. We
report it for completeness with two observations that survive even at this size.

First, the coarse \texttt{condition} column in the artifact pools
\emph{both} injected faults into a single 18-row bucket; grouping on it produces a
``fault'' row that is really two different faults averaged together, which we
avoid by breaking rows out by fault. Second, plan-and-execute succeeds in
\ArchPlanCtl{} of \ArchCellN{} control trials against \ArchReactCtl{} for ReAct,
while taking \ArchPlanTime{}~s against \ArchReactTime{}~s --- a $4.7\times$
wall-clock difference. A comparison at equal \emph{step} budgets is therefore not
a comparison at equal \emph{time} budgets, and the artifact records no per-role
call counters, so the two arms cannot be shown to have had equal LLM budgets. No
architecture ranking follows from this, in either direction.

\subsection{The recovery artifact is degenerate}
\label{sec:res-recovery}

The recovery artifact is the paper's second methodological warning, with the
column-by-column detail in \appref{app:recovery}. It records behavioural flags ---
whether the observation was refreshed, whether the tool changed, whether a stale
action was repeated --- which is exactly the process-level evidence one would want
for attribution. In practice \RecDegenerateCount{} of its columns are constant
across all \RecN{} rows and carry no information, \texttt{plan\_revised} was never
recorded at all, and \texttt{recovered} is row-for-row identical to
\texttt{official\_success}: not an independent measurement of recovery, but a
restatement of the outcome it would be used to explain.

We report it rather than silently dropping it because the failure mode is
instructive: a recovery-behaviour artifact can look rich --- a dozen behavioural
columns --- while containing two varying quantities, one of which is the outcome.
Process-level attribution needs columns that vary \emph{independently} of the
outcome, and that has to be checked before the analysis, not after.
